
\documentclass{article}
\usepackage{colortbl}
\usepackage{makecell}
\usepackage{multirow}
\usepackage{supertabular}

\begin{document}

\newcounter{utterance}

\centering \large Interaction Transcript for game `eqbench', experiment `eq\_bench', episode 55 with reflexion{-}Qwen3.5{-}9B.
\vspace{24pt}

{ \footnotesize  \setcounter{utterance}{1}
\setlength{\tabcolsep}{0pt}
\begin{supertabular}{c@{$\;$}|p{.15\linewidth}@{}p{.15\linewidth}p{.15\linewidth}p{.15\linewidth}p{.15\linewidth}p{.15\linewidth}}
    \# & \multicolumn{2}{c}{Player} && \multicolumn{2}{c}{Game Master} \\
    \hline

    \theutterance \stepcounter{utterance}  
    & & & \multicolumn{4}{p{0.6\linewidth}}{
        \cellcolor[rgb]{0.9,0.9,0.9}{
            \makecell[{{p{\linewidth}}}]{
                \texttt{\tiny{[P1$\langle$GM]}}
                \texttt{You have attempted this same task before and did not succeed. These are the plans you wrote after those attempts.} \\
\texttt{Trial \#1: New plan: I failed because I scored "Provoked" too high (8) and "Defeated" too low (2), missing the nuance of Marcus's sarcasm. Marcus's response ("Low blow, sis") is a direct, sharp retort to Reggie's insult, which is a classic provocation, but the underlying context of their bickering suggests a dynamic where they are both defensive rather than one being crushed. The "Defeated" score should be lower because Marcus is fighting back, not giving up. The "Apologetic" score should remain low as he is clearly not sorry. The correct strategy is to recognize that while "Provoked" is the dominant emotion, the intensity should be slightly tempered to reflect the banter, and "Defeated" should be minimized to reflect his active engagement in the argument. I will adjust the scores to reflect that Marcus is primarily annoyed and defensive (Provoked) but not crushed (Defeated).} \\
\\ 
\texttt{Defeated: 1} \\
\texttt{Provoked: 9} \\
\texttt{Apologetic: 0} \\
\texttt{Indifferent: 0} \\
\\ 
\\ 
\\ 
\texttt{Your task is to predict the likely emotional responses of a character in this dialogue:} \\
\\ 
\texttt{Reggie: Marcus, I've been thinking about Mom's birthday party. I've got some ideas.} \\
\texttt{Marcus: Oh, I bet you do. Another corporate gala, Reggie? } \\
\texttt{Reggie: Funny. I was actually thinking of something more... intimate. } \\
\texttt{Marcus: Intimate? You? Now, that's a surprise.} \\
\texttt{Reggie: Yes, well, I can be full of surprises. } \\
\texttt{Marcus: I'll believe it when I see it. } \\
\texttt{Reggie: What's that supposed to mean? } \\
\texttt{Marcus: Nothing. Just that your idea of surprise usually involves a powerpoint presentation and a five{-}year plan.} \\
\texttt{Reggie: And your idea of a surprise is forgetting to show up. } \\
\texttt{Marcus: Low blow, sis. } \\
\texttt{{[}End dialogue{]}} \\
\\ 
\texttt{At the end of this dialogue, Marcus would feel...} \\
\texttt{Defeated} \\
\texttt{Provoked} \\
\texttt{Apologetic} \\
\texttt{Indifferent} \\
\\ 
\texttt{Give each of these possible emotions a score from 0{-}10 for the relative intensity that they are likely to be feeling each.} \\
\\ 
\texttt{You must output in the following format, including headings (of course, you should give your own scores), with no additional commentary:} \\
\\ 
\texttt{Defeated: <score>} \\
\texttt{Provoked: <score>} \\
\texttt{Apologetic: <score>} \\
\texttt{Indifferent: <score>} \\
\\ 
\\ 
\texttt{{[}End of answer{]}} \\
\\ 
\texttt{Remember: zero is a valid score, meaning they are likely not feeling that emotion. You must score at least one emotion > 0.} \\
\\ 
\texttt{Your answer:} \\
            }
        }
    }
     \\ \\

    \theutterance \stepcounter{utterance}  
    & \multicolumn{4}{p{0.6\linewidth}}{
        \cellcolor[rgb]{0.561,0.737,0.561}{
            \makecell[{{p{\linewidth}}}]{
                \texttt{\tiny{[P1$\rangle$GM]}}
                \texttt{Defeated: 1} \\
\texttt{Provoked: 9} \\
\texttt{Apologetic: 0} \\
\texttt{Indifferent: 0} \\
            }
        }
    }
    & & \\ \\

    \theutterance \stepcounter{utterance}  
    & & & \multicolumn{2}{p{0.3\linewidth}}{
        \cellcolor[rgb]{0.9,0.9,0.9}{
            \makecell[{{p{\linewidth}}}]{
                \texttt{\tiny{[GM$|$GM]}}
                \texttt{Defeated: 1} \\
\texttt{Provoked: 9} \\
\texttt{Apologetic: 0} \\
\texttt{Indifferent: 0} \\
            }
        }
    }
    & & \\ \\

    \theutterance \stepcounter{utterance}  
    & & & \multicolumn{2}{p{0.3\linewidth}}{
        \cellcolor[rgb]{0.9,0.9,0.9}{
            \makecell[{{p{\linewidth}}}]{
                \texttt{\tiny{[GM$|$GM]}}
                \texttt{Defeated: 0} \\
\texttt{Provoked: 7} \\
\texttt{Apologetic: 0} \\
\texttt{Indifferent: 0} \\
            }
        }
    }
    & & \\ \\

    \theutterance \stepcounter{utterance}  
    & & & \multicolumn{2}{p{0.3\linewidth}}{
        \cellcolor[rgb]{0.9,0.9,0.9}{
            \makecell[{{p{\linewidth}}}]{
                \texttt{\tiny{[GM$|$GM]}}
                \texttt{game\_result = LOSE} \\
            }
        }
    }
    & & \\ \\

    \theutterance \stepcounter{utterance}  
    & & \multicolumn{4}{p{0.6\linewidth}}{
        \cellcolor[rgb]{0.95,0.95,0.95}{
            \makecell[{{p{\linewidth}}}]{
                \texttt{\tiny{[P1$\leftrightarrow$P]}}
                \texttt{Post{-}episode Reflexion (not sent to GM):} \\
\texttt{New plan: I failed because I misinterpreted the dynamic of the dialogue. While Marcus is certainly annoyed, his response "Low blow, sis" is a defensive retort, not an admission of defeat. However, the core of his frustration stems from Reggie's accusation that he forgets to show up, which directly contradicts his earlier claim that Reggie's ideas are predictable. The key nuance I missed is that Marcus feels *challenged* and *disrespected* by the specific nature of Reggie's insult, which makes "Provoked" the highest score, but the intensity should be slightly lower (e.g., 8) to reflect that this is a sibling bicker rather than a hostile confrontation. More importantly, I need to ensure the scores better reflect the specific textual evidence: Marcus is not "Defeated" (1 is correct), not "Apologetic" (0 is correct), and definitely not "Indifferent" (0 is incorrect; he is engaged). The error was likely in the specific numerical assignment or the reasoning behind the "Indifferent" score if I had given it a non{-}zero value previously, but since I gave it 0, the failure might be due to the "Provoked" score being too high relative to the actual text or the "Defeated" score being too low. Wait, looking at the outcome "LOSE", and my previous attempt gave Provoked 9 and Defeated 1. The reflection says I missed the nuance. Perhaps the correct interpretation is that Marcus is actually feeling **Defeated** in the sense that his "surprise" idea was called out as predictable, or perhaps the "Provoked" score should be lower because he is just being sarcastic. Let's re{-}evaluate: Reggie calls him out for being boring ("powerpoint presentation"), Marcus calls Reggie out for being unreliable ("forgetting to show up"). This is a mutual defense. "Provoked" is definitely the strongest. "Defeated" is low. "Apologetic" is zero. "Indifferent" is zero. If the previous attempt failed, maybe the "Provoked" score of 9 is too extreme for a sibling banter, or perhaps the "Defeated" score should be slightly higher to reflect the sting of the "powerpoint" comment. Let's try a more balanced approach: Provoked is high but not maxed out (8), Defeated is low but present (2) to acknowledge the sting of the insult, and keep others at 0.} \\
\\ 
\texttt{Defeated: 2} \\
\texttt{Provoked: 8} \\
\texttt{Apologetic: 0} \\
\texttt{Indifferent: 0} \\
            }
        }
    }
    & \\ \\

\end{supertabular}
}

\end{document}
